\chapter{Estado da Arte}\label{cap:estadoarte}
\addcontentsline{toc}{chapter}{Estado da Arte}

Este capítulo consolida os trabalhos recentes diretamente relacionados ao problema deste TCC: uma comparação empírica, sob protocolo único e em aplicações Python, entre SAST tradicional, ferramentas empresariais de desenvolvimento assistido por IA e uma abordagem híbrida SAST + IA. Cada seção isola uma peça desse desenho experimental e termina apontando o que falta para viabilizá-lo. 

A Seção ~\ref{sec:avaliacoes-sast} mostra como ferramentas SAST têm sido avaliadas e por que seu comportamento real é o ponto de partida (e o baseline) do experimento. 

A Seção ~\ref{sec:llms-deteccao-vulnerabilidades} mostra como LLMs isolados vêm sendo aplicados à mesma tarefa, distinguindo-os das ferramentas-agente avaliadas neste TCC. 

A Seção ~\ref{sec:comparacoes-sast-ia} revisa os estudos que já colocam SAST e LLMs lado a lado , a linha de pesquisa mais próxima deste trabalho  e expõe o quanto ela ainda ignora Python e ferramentas comerciais. 

A Seção ~\ref{sec:abordagens-hibridas}  trata do terceiro braço, o híbrido, mapeando os padrões de combinação já testados. 
s
A Seção ~\ref{sec:agentes-ia-seguranca} examina o paradigma de agentes de IA em segurança de código, do qual emergem Cursor e OpenAI Codex; Claude Code é tratado apenas como alternativa comercial contextual. 

A Seção ~\ref{sec:lacunas-literatura} sintetiza as lacunas remanescentes — inclusive o resultado de uma busca complementar dedicada, especificamente para verificar se a lacuna mais crítica identificada neste levantamento já havia sido preenchida pela literatura mais recente.


\section{Avaliações recentes de ferramentas SAST}\label{sec:avaliacoes-sast}

Antes de comparar SAST com qualquer abordagem baseada em IA, é preciso estabelecer o que o SAST tradicional efetivamente entrega quando avaliado com rigor , esse é o baseline contra o qual os braços de ferramentas-agente empresariais e híbrido deste TCC serão medidos, e por isso a literatura revisada nesta seção é fundação direta do desenho experimental. As avaliações empíricas recentes convergem em um diagnóstico: ferramentas SAST tendem a ser precisas quando alertam, mas deixam de detectar a maior parte das vulnerabilidades reais. Em benchmark sintético de larga escala, a melhor entre oito ferramentas atingiu \textit{recall} médio de aproximadamente 0,58 mesmo nas categorias que declarava cobrir \cite{Esposito2024SASTComparison}. Em código real, o quadro é mais severo: sete ferramentas gratuitas para Java, que atingiam F1 de até 84\% em benchmark sintético, detectaram no máximo 12,7\% de 165 CVEs reais, e mesmo a combinação de quatro ferramentas deixou 70,9\% dos casos sem detecção \cite{Li2023SASTJava}. Avaliação de quatro ferramentas em 170 \textit{commits} com CVEs confirmou taxas individuais entre 11\% e 27\%, com a combinação alcançando 38,8\% \cite{Bennett2024Semgrep}, e um estudo com nove ferramentas para JavaScript observou taxa média de verdadeiros positivos de 15,1\%, com precisão extremamente baixa nas ferramentas de maior cobertura \cite{Brito2023JSSAST}. No maior estudo longitudinal disponível, o CodeQL detectou apenas 4,4\% de quase 4 mil CVEs analisáveis, incluindo apenas 36 casos em Python, além de apresentar regressões de detecção entre versões \cite{Noirot2025CodeQLLongitudinal}.

Dois fatores explicativos se repetem. O primeiro é a dependência do conjunto de regras: a taxa de detecção do Semgrep quase triplicou (15,3\% para 44,7\%) após engenharia manual de regras \cite{Bennett2024Semgrep}, e a dissecção de falsos negativos em ferramentas baseadas em consultas mostrou que a maioria decorre de capacidade inadequada das regras, e não do motor de análise \cite{Li2024SASTCC, Li2023SASTJava}. Em Python, o mesmo padrão foi observado: Bandit e Semgrep em configuração padrão detectaram, respectivamente, cerca de 59,7\% e 28,6\% de 108 vulnerabilidades reais de pacotes PyPI, valores que subiram para 72–81\% com regras customizadas construídas ao custo de meses de trabalho especializado \cite{Marquer2026SAGA}. O segundo fator é a complementaridade entre ferramentas: nenhuma domina todas as categorias, e a combinação eleva a cobertura de forma consistente \cite{Esposito2024SASTComparison, Zhu2024SASTAndroid, Dalaq2025SASTSLR}.

Paralelamente, a linha de pesquisa sobre adoção reforça que o custo de triagem é o gargalo prático: falsos positivos são a barreira mais citada na literatura \cite{Wadhams2024SASTBarriers}, grande parte dos alertas permanece anos sem tratamento \cite{Aloraini2019SASTWarnings}, e profissionais da indústria demandam diagnósticos mais finos do que métricas agregadas \cite{Li2025IndustrySAST}. Em conjunto, esses estudos estabelecem o perfil do braço SAST do presente experimento determinístico, rápido e barato, porém com cobertura limitada e relatórios de baixa acionabilidade e indicam as dimensões em que abordagens baseadas em IA poderiam agregar valor. É justamente esse potencial de complementaridade que a seção seguinte examina, do lado oposto do espectro: o que os LLMs, isoladamente, conseguem detectar quando confrontados com a mesma tarefa.


\section{LLMs para detecção de vulnerabilidades em código}\label{sec:llms-deteccao-vulnerabilidades}
Se a Seção Seção~\ref{sec:avaliacoes-sast} caracterizou o braço SAST do experimento, esta seção caracteriza o material bruto do segundo braço, sendo a capacidade de LLMs de propósito geral em identificar vulnerabilidades , que mais adiante será encapsulado, no caso deste TCC, pelas ferramentas empresariais de desenvolvimento assistido por IA (Seção ~\ref{sec:agentes-ia-seguranca}). Entender onde os LLMs falham quando avaliados diretamente é pré-condição para interpretar corretamente o desempenho dessas ferramentas quando operam como agentes sobre um repositório. A literatura sobre LLMs para detecção de vulnerabilidades cresceu rapidamente: revisões sistemáticas recentes catalogam de 58 a 208 estudos primários \cite{Sheng2025LLMSurvey, Zhou2025LLMReview, Germano2025LLMSLR}. O quadro geral é de potencial acompanhado de fragilidades. No maior estudo empírico com modelos de propósito geral, 16 LLMs avaliados sobre 5 mil trechos de código atingiram acurácia média de 62,8\% e F1 médio de 0,71, com queda de cerca de 10 pontos percentuais ao passar de dados sintéticos para dados reais \cite{Khare2025LLMEffectiveness}. Avaliações mais rigorosas são ainda menos otimistas: quando benchmarks tradicionais foram corrigidos quanto a rótulos errados e vazamento de dados, o F1 de modelos de código despencou de 68,3\% para 3,1\%, e nem o GPT-4 com raciocínio em cadeia superou o acaso na distinção entre função vulnerável e sua versão corrigida \cite{Ding2025CodeLMVuln}. Estudos de raciocínio confirmam o diagnóstico: LLMs exibem consistência semântica inferior a 50\% em tarefas estruturadas de análise de vulnerabilidades \cite{Li2025SVTrustEval}, erram a classificação da CWE específica mesmo quando "percebem" o problema (cerca de 16\% de acerto, próximo do aleatório) \cite{Gaspar2025OpenLLM, Lin2025LLMCompare} e degradam a localização em arquivos longos \cite{Sovrano2025LostInEnd}.

Um segundo eixo de resultados diz respeito ao papel do contexto. Avaliações com trechos isolados de código subestimam ou distorcem a capacidade dos modelos; com contexto adequado do repositório, a acurácia estrita chegou a 67\% no benchmark CORRECT, e a maioria dos falsos positivos foi atribuída a erros de raciocínio, não de classificação \cite{Li2025LLMVDAsk}. Na escala de projeto, porém, detectores baseados em LLM apresentaram \textit{recall} entre 21\% e 34\% com taxas de falsos alertas de 85–97\%, custo computacional elevado e causas de erro dominadas por análise interprocedural rasa \cite{Li2026ProjectScaleLLM}. Abordagens que fornecem contexto estruturado ao modelo  fatias de código guiadas por grafos de propriedade \cite{Lekssays2025LLMxCPG}, abstrações de funções chamadas \cite{Yang2025ContextVD} ou especificações extraídas de correções históricas \cite{Zhu2025SpecGuidedVD} , melhoram consistentemente os resultados, o que sugere que a forma de entregar o código ao modelo importa tanto quanto o modelo em si. Essa observação é diretamente relevante para o TCC: Cursor e Codex operam com camadas de contexto e ferramentas sobre o repositório \cite{Cursor2025Docs, OpenAI2025CodexDocs}; isso os distingue dos LLMs isolados revisados nesta seção, sem permitir presumir a priori superioridade de desempenho.

Uma terceira vertente investiga modelos especializados por \textit{fine-tuning}. Para Python especificamente, o SecureQwen reportou F1 entre 84\% e 99\% por categoria CWE, embora com rótulos gerados por um \textit{ensemble} de ferramentas SAST (incluindo Bandit e Semgrep), o que introduz circularidade na avaliação \cite{Mechri2025SecureQwen}. Resultados similares de alto desempenho em benchmarks controlados foram reportados para C/C++ \cite{Ferrag2025SecureFalcon}, mas os ganhos do \textit{fine-tuning} sobre linhas de base são frequentemente modestos em dados desbalanceados \cite{Shestov2025LLMFinetune}. Cabe notar que esses modelos dedicados diferem substancialmente das ferramentas empresariais de desenvolvimento assistido por IA avaliadas neste trabalho, que utilizam modelos generalistas com acesso a contexto de repositório. Por fim, a dimensão da explicação permanece pouco explorada: apenas 5,3\% dos 208 estudos de uma revisão sistemática abordam explicação de vulnerabilidades, contra 91,3\% dedicados à detecção binária \cite{Germano2025LLMSLR}; o trabalho que propôs rubricas de corretude, clareza e acionabilidade constatou que alta taxa de detecção não implica explicações acionáveis \cite{Mao2025ExplainableVD}. Tomados em conjunto, os dois primeiros eixos revisados nesta seção , desempenho isolado modesto e sensibilidade forte ao contexto motivam diretamente a pergunta central da próxima seção: quando SAST e LLM são colocados lado a lado sob o mesmo protocolo, qual abordagem efetivamente prevalece, e sob quais condições?


\section{Comparações entre SAST tradicional e abordagens baseadas em IA}\label{sec:comparacoes-sast-ia}
As comparações diretas entre SAST e LLMs formam a linha de pesquisa mais próxima deste TCC: é aqui que o \textit{trade-off} entre os dois braços do experimento é medido explicitamente, sob o mesmo conjunto de dados. O estudo pioneiro comparou o ChatGPT com onze ferramentas SAST para PHP e encontrou o padrão que se repetiria nos trabalhos seguintes: o LLM detectou 62–68\% das vulnerabilidades, contra 32\% da melhor ferramenta tradicional, mas alertou incorretamente 91\% dos trechos seguros e variou 6 pontos percentuais entre execuções idênticas \cite{Ozturk2023AIvsSAST}. Avaliação posterior com cinco LLMs no mesmo domínio confirmou taxas de detecção de até 61,5\% acompanhadas de falsos positivos superiores a 60\% em código corrigido \cite{Cetin2024LLMPHP}. Estudos anteriores com modelos da geração GPT-3 já apontavam a mesma assimetria ,\textit{recall} elevado com taxa de falsos positivos várias vezes superior à das ferramentas estáticas \cite{Purba2023LLMVuln}.

Comparações metodologicamente mais rigorosas refinaram o quadro. Em benchmark controlado sobre 578 arquivos Java, o melhor LLM naquele momento (GPT-4 com raciocínio em cadeia, prompt CoT de 8 passos) atingiu F1 de 0,672 e cobriu 13 das 17 categorias CWE, contra F1 de 0,655 e cobertura de 11 categorias da melhor ferramenta SAST; em contrapartida, o SAST produziu menos falsos positivos (23 contra 55), com custo financeiro nulo e execução em menos de um minuto \cite{Tamberg2025LLMHarness}. Em dez projetos C\# reais, LLMs de fronteira superaram três ferramentas SAST em F1 (0,75–0,80 contra 0,26–0,55), mas falharam completamente na localização precisa (linha/coluna) dos achados \cite{Gnieciak2025LLMvsSAST}. Um estudo de larga escala com 15 ferramentas SAST e 12 LLMs de código aberto sobre repositórios reais em Java, C e Python — um dos poucos a incluir Python nesta escala , reportou o mesmo padrão qualitativo já visto acima (SAST com baixa cobertura e poucos falsos positivos; LLMs com cobertura muito superior e falsos positivos elevados, mitigados apenas parcialmente por \textit{ensemble}), embora se trate de preprint ainda sem revisão por pares quase dois anos após a submissão, o que recomenda cautela redobrada quanto aos valores absolutos reportados \cite{Zhou2024SASTLLMRepo}.

No único benchmark identificado com foco em Python e revisão por pares pendente, cobrindo 26 repositórios e 796 entradas rotuladas, os \textit{scanners} baseados em LLMs de propósito geral superaram amplamente as ferramentas de regras , o Semgrep obteve F3 de 17,7, contra 51,7 de um LLM comercial e 73,0 da melhor ferramenta especializada \cite{Pellew2026RealVuln}. Também em Python, um estudo comparando aprendizado de máquina, ChatGPT e análise estática sobre 21 scripts concluiu que as abordagens baseadas em IA superaram significativamente as ferramentas estáticas, embora a amostra reduzida limite a generalização \cite{Farasat2025PythonSecurity}. Um preprint mais recente (janeiro de 2026) propõe um protocolo de avaliação hierárquica por CWE especificamente para Bandit, Semgrep e CodeQL contra LLMs abertos em nível de função em Python, reforçando que Python começa a receber atenção comparativa dedicada, ainda que os resultados numéricos desse trabalho não tenham podido ser verificados de forma independente nesta revisão e, por isso, não sejam aqui reproduzidos \cite{Adnan2026ALPHA}.


\begin{table}[htbp]
\centering
\caption{Estudos comparativos SAST–LLM mais próximos do presente trabalho.}\label{tab:compartivos-sast-llm}
\begin{tabularx}{\textwidth}{|X|X|X|X|X|}
\hline
\textbf{Estudo} & \textbf{Linguagem e dados} & \textbf{Comparação} & \textbf{Métricas e achados selecionados} & \textbf{Situação} \\
\hline
Ozturk et al. \cite{Ozturk2023AIvsSAST} & PHP & ChatGPT vs. 11 SAST & Detecção de 62–68\% vs. 32\%; 91\% de falsos alertas em trechos seguros; variação de 6 p.p. & Estudo pioneiro \\
\hline
Tamberg et al. \cite{Tamberg2025LLMHarness} & Java, 578 arquivos & GPT-4 com raciocínio em cadeia vs. SAST & F1 de 0,672 vs. 0,655; 13 vs. 11 CWEs; 55 vs. 23 falsos positivos & Revisado por pares \\
\hline
Gnieciak e Szandala \cite{Gnieciak2025LLMvsSAST} & C\#, 10 projetos & LLMs de fronteira vs. 3 SAST & F1 de 0,75–0,80 vs. 0,26–0,55; falha de localização precisa pelos LLMs & Preprint \\
\hline
Zhou et al. \cite{Zhou2024SASTLLMRepo} & Java, C e Python; repositórios & 15 SAST vs. 12 LLMs & Maior cobertura dos LLMs, com maior proporção de funções marcadas & Preprint \\
\hline
Pellew e Raza \cite{Pellew2026RealVuln} & Python, 26 repositórios, 796 entradas & Ferramentas de regras vs. \textit{scanners} LLM & F3: Semgrep 17,7; LLM comercial 51,7; melhor ferramenta especializada 73,0 & Preprint \\
\hline
Farasat e Posegga \cite{Farasat2025PythonSecurity} & Python, 21 scripts & Aprendizado de máquina, ChatGPT e SAST & Abordagens de IA superaram SAST; amostra reduzida & Revisado por pares \\
\hline
\end{tabularx}f
\end{table}

A Tabela ~\ref{tab:compartivos-sast-llm} torna explícita a heterogeneidade da evidência: os trabalhos diferem em linguagem, granularidade, conjunto de dados e métrica. Por essa razão, não é metodologicamente apropriado agregar seus números em um gráfico único; a tabela preserva a comparabilidade limitada e evita inferências que os estudos não sustentam.

Em síntese, a literatura comparativa sugere um \textit{trade-off} consistente: LLMs oferecem maior cobertura e explicações em linguagem natural, ao custo de falsos positivos elevados, não-determinismo, custo financeiro e localização imprecisa; ferramentas SAST oferecem precisão, reprodutibilidade e custo baixo, ao custo de cobertura limitada \cite{Ozturk2023AIvsSAST, Tamberg2025LLMHarness, Gnieciak2025LLMvsSAST}. Nenhum desses estudos, contudo, avaliou ferramentas empresariais de desenvolvimento assistido por IA operando como agentes sobre o repositório ,cenário deste TCC , e mesmo somando os trabalhos mais recentes localizados nesta revisão, o ecossistema Python permanece coberto majoritariamente por preprints \cite{Pellew2026RealVuln, Farasat2025PythonSecurity, Zhou2024SASTLLMRepo, Adnan2026ALPHA}. Diante de um \textit{trade-off} tão bem documentado entre as duas abordagens, é natural que uma parcela crescente da literatura tenha deixado de tratá-las como alternativas excludentes e passado a investigar formas de combiná-las , é o que a seção seguinte revisa.




\section{Abordagens híbridas SAST + IA}\label{sec:abordagens-hibridas}
Diante do \textit{trade-off} descrito na seção anterior, uma linha crescente de trabalhos combina análise estática e LLMs , o terceiro braço deste TCC, e o que a literatura sugere ser o de maior potencial prático, ainda que o mais carente de validação padronizada. Três padrões arquiteturais predominam. No primeiro o mais próximo do desenho deste TCC , o LLM atua a jusante do SAST, validando, filtrando e explicando os alertas. Um sistema multi-agente que revisa candidatos do Semgrep com contexto do código reduziu a taxa de falsos positivos de 0,423 para 0,048 (redução de 88,6\%) ao custo de apenas 3,1\% dos verdadeiros positivos, elevando o F1 de 0,784 para 0,909 no OWASP Benchmark \cite{Ameen2026QASecClaw}. Resultados da mesma ordem foram reportados ao tratar saídas SAST como contratos estruturados antes da adjudicação por LLM (F1 de até 0,955 em projetos reais) \cite{Iranmanesh2025ZeroFalse}, na triagem de relatórios com raciocínio LLM (precisão de 9,4\% para 52,1\%) \cite{Li2025SASTReportsLLM} e na integração Semgrep + LLM com redução de aproximadamente 91\% dos alertas \cite{Agrawal2025SASTGenius}. Na comparação experimental de cenários, o híbrido SAST + LLM superou ambas as abordagens isoladas nos cenários de menor criticidade, embora o SAST puro tenha permanecido preferível quando o \textit{recall} é prioritário, e o híbrido tenha inclusive piorado o desempenho em categorias específicas \cite{Ciavatta2026LLMSAST}.

No segundo padrão, o LLM atua a montante, ampliando a capacidade do motor estático: inferindo especificações de \textit{taint} (\textit{sources} e \textit{sinks}) para o CodeQL, o que dobrou a detecção de CVEs reais (de 22,5\% para 45,8\%) \cite{Li2025IRIS}; descobrindo \textit{sources/sinks/sanitizers} em escala industrial \cite{Luo2025SASTLLMDFA}; gerando consultas CodeQL a partir de descrições de CWE \cite{Wang2026CQLLM}; ou sintetizando regras de forma autônoma com validação por prova de conceito, abordagem que descobriu CVEs novos em pacotes Python do PyPI \cite{Tsigkourakos2026QRS}. O terceiro padrão emprega o SAST como oráculo de realimentação em laços de geração e correção: em Python, o retorno de alertas do Bandit ao LLM reduziu o código vulnerável gerado de cerca de 40\% para menos de 10\% \cite{Alrashedy2025FDSP, Lee2026CodeEnhancer}. Um estudo independente, em escala de produção (733 trechos de código Python e JavaScript gerados por GitHub Copilot, CodeWhisperer e Codeium extraídos de projetos reais do GitHub, com 29,5\% dos trechos Python contendo fraquezas de segurança em 43 categorias de CWE), confirma esse mesmo padrão de realimentação SAST→LLM fora de ambiente controlado: fornecer ao Copilot Chat os alertas gerados por ferramentas de análise estática permitiu corrigir até 55,5\% das fraquezas identificadas — evidência revisada por pares, em Python, de que uma ferramenta comercial de IA pode atuar como agente corretor quando municiada por sinal de SAST, mesmo que o artigo não avalie a mesma ferramenta como agente \textit{detector} \cite{Fu2025CopilotSecurity}.

A leitura crítica desses resultados exige três ressalvas. Primeiro, boa parte dos trabalhos híbridos mais próximos do TCC são preprints ainda não revisados por pares \cite{Ameen2026QASecClaw, Iranmanesh2025ZeroFalse, Agrawal2025SASTGenius, Tsigkourakos2026QRS}. Segundo, os experimentos concentram-se em Java (frequentemente no OWASP Benchmark sintético) e Solidity, com escassa validação em Python \cite{Ciavatta2026LLMSAST, Ameen2026QASecClaw, Li2025SASTReportsLLM} — o estudo de Fu et al. é uma das poucas exceções revisadas por pares com dados reais de Python, ainda que restrito ao padrão de correção, não de detecção \cite{Fu2025CopilotSecurity}. Terceiro, a combinação não é garantidamente benéfica: a união ingênua de detectores baseados em regras e em aprendizado piora o desempenho global \cite{Croft2021RuleLearningSAST}, e a filtragem por LLM pode degradar categorias específicas de vulnerabilidade \cite{Ciavatta2026LLMSAST, Ameen2026QASecClaw}. Surveys dedicados à integração LLM + análise estática convergem, ainda assim, para a conclusão de que a combinação supera cada abordagem isolada na maioria dos cenários, posicionando o paradigma híbrido como direção promissora, porém carente de validação empírica padronizada \cite{Jin2025LLMSASTSurvey, Khare2025LLMEffectiveness}. Falta a essa linha de pesquisa, contudo, um ingrediente específico: os padrões híbridos revisados aqui usam LLMs via API, orquestrados por \textit{pipelines} construídos pelos próprios autores, nenhum emprega diretamente uma ferramenta empresarial de desenvolvimento assistido por IA operando de forma autônoma e agêntica sobre um repositório, o que remete à seção seguinte.


\section{Agentes de IA e workflows automatizados para segurança de código}\label{sec:agentes-ia-seguranca}
Esta seção trata do paradigma do qual emergem as ferramentas avaliadas neste TCC, sendo eles Cursor e OpenAI Codex: agentes de IA comerciais que operam em ciclos de raciocínio e ação sobre um repositório, com acesso a ferramentas. Claude Code é mencionado somente como alternativa comercial contextual, não como objeto de avaliação. Entender o que a literatura acadêmica já demonstrou (e o que ainda não demonstrou) sobre esse paradigma aplicado à segurança de código é o elo final antes de justificar formalmente o desenho experimental na Seção ~\ref{sec:lacunas-literatura}. O paradigma foi consolidado em engenharia de software por benchmarks de resolução de \textit{issues} reais \cite{Jimenez2024SWEbench} e pela demonstração de que interfaces agente-computador especializadas multiplicam o desempenho dos modelos \cite{Yang2024SWEagent}. Cursor e Codex implementam esse paradigma com modos de agente, acesso ao repositório e ferramentas de execução configuráveis \cite{Cursor2025Docs, OpenAI2025CodexDocs}. Sua transposição para segurança de código, porém, é recente e apresenta resultados mistos. No benchmark JITVUL, agentes do tipo ReAct com acesso a ferramentas melhoraram a acurácia em pares em relação ao \textit{prompt} simples em alguns cenários, mas o desempenho absoluto permaneceu baixo (inferior a 21\%) e o simples enriquecimento do \textit{prompt} com contexto de dependências por vezes superou o agente completo \cite{Yildiz2025LLMAgents}.

Arquiteturas multi-agente mais elaboradas mostram resultados promissores em cenários específicos. Um \textit{ensemble} de agentes orquestrando análise estática com recuperação de conhecimento de vulnerabilidades encontrou falhas em bases de código industriais nas quais o CodeQL isolado nada detectou, com custo da ordem de poucos dólares por repositório \cite{Liang2026Argus}; a tríade neuro-simbólica de agentes que sintetiza e valida regras descobriu cinco CVEs novos em pacotes populares do PyPI \cite{Tsigkourakos2026QRS}; e agentes de filtragem sustentam os melhores resultados de redução de falsos positivos \cite{Ameen2026QASecClaw}. Um exemplo particularmente informativo, publicado ems veículo revisado por pares de primeira linha (ICSE 2026), é o VulTrial: uma arquitetura multi-agente inspirada em um julgamento (agentes com papéis de pesquisador de segurança, autor do código, moderador e conselho de revisão) que, sobre o benchmark PrimeVul-pair (C/C++), quase dobrou o número de pares corretamente classificados em relação a \textit{baselines} single-agent e multi-agente simples (81 contra 40–44 pares corretos com GPT-4o, chegando a 96 com ajuste do agente moderador) e identificou três vulnerabilidades zero-day confirmadas em projetos reais em produção e reportando,  o custo em dólares de cada configuração (de US\$ 7,46 com GPT-3.5 a cerca de US\$ 23,90 com GPT-4o para o conjunto de teste completo), uma das poucas evidências quantitativas de custo disponíveis nesta linha de pesquisa \cite{Widyasari2026VulTrial}. Ainda assim, o VulTrial usa LLMs via API orquestrados por uma arquitetura de agentes construída pelos próprios autores , não Cursor, Codex ou Claude Code e é validado exclusivamente em C/C++, não em Python.

No extremo oposto da cadeia, o único estudo de campo com ferramenta de detecção e reparo integrada ao IDE encontrou utilidade percebida baixa (2,5/5 para alertas) e 75\% de correções não aplicáveis, evidenciando a distância entre benchmark e prática \cite{Steenhoek2025IDEUserStudy}. Uma busca complementar dedicada especificamente a essa lacuna localizou um preprint que avalia o recurso "Code Review" do GitHub Copilot; seu achado qualitativo é que a ferramenta falha em identificar vulnerabilidades críticas e concentra comentários em problemas de estilo \cite{Amro2025CopilotReview}. Há também evidência acadêmica direta, porém adjacente, para Cursor: o SecCodePLT avaliou manualmente 153 exemplos Python de geração e completamento seguro e observou taxas de código seguro de 52,8\% em testes dinâmicos e 62,0\% em testes baseados em regras, sem política de segurança; com política explícita, os valores subiram para 67,4\% e 86,7\%, respectivamente \cite{Nie2025SecCodePLT}. Esse estudo avalia a segurança do código que Cursor gera, e não sua capacidade de detectar vulnerabilidades já existentes em um repositório. A Tabela ~\ref{tab:ferramentas-empresariais-estudo} delimita essas evidências próximas e evita confundir geração segura, revisão de \textit{pull request} e detecção de vulnerabilidades.



\begin{table}[htbp]
\centering
\caption{Ferramentas empresariais de IA em estudos de segurança adjacentes ao TCC}\label{tab:ferramentas-empresariais-estudo}
\begin{tabularx}{\textwidth}{|X|X|X|X|X|}
\hline
\textbf{Ferramenta ou sistema} & \textbf{Estudo} & \textbf{Cenário} & \textbf{Métrica ou achado} & \textbf{Relação com o TCC} \\
\hline
Cursor & Nie et al. \cite{Nie2025SecCodePLT} & 153 exemplos Python; geração e completamento seguro & 52,8\% em testes dinâmicos e 62,0\% em testes de regras; políticas de segurança elevaram os valores para 67,4\% e 86,7\% & Avaliação direta de Cursor em segurança, mas não de detecção \\
\hline
GitHub Copilot Code Review & Amro e Alalfi \cite{Amro2025CopilotReview} & Revisão de código vulnerável & Falhas qualitativas em vulnerabilidades críticas; sem métricas verificadas nesta revisão & Produto comercial adjacente; preprint \\
\hline
Protótipo integrado ao IDE & Steenhoek et al. \cite{Steenhoek2025IDEUserStudy} & Estudo de campo de detecção e reparo & Utilidade dos alertas: 2,5/5; 75\% das correções não aplicáveis & Não avalia Cursor ou Codex \\
\hline
Claude Code & Luo et al. \cite{Luo2026CodeAugur} & \textit{Baseline} em benchmark de agentes com validação dinâmica & Comparado a um agente acadêmico de inferência de especificações e \textit{fuzzing} & Alternativa comercial contextual; preprint e fora de Python \\
\hline
\end{tabularx}
\end{table}

Portanto, as evidências encontradas não avaliam Cursor ou OpenAI Codex, configurados como ferramentas-agente de desenvolvimento, para detectar vulnerabilidades em aplicações Python. Elas reforçam a relevância do tema, mas tratam de geração segura, revisão de código, protótipos de IDE ou agentes acadêmicos em outros contextos.


\section{Lacunas identificadas na literatura}\label{sec:lacunas-literatura}
A análise consolidada nas seções anteriores permite identificar as seguintes lacunas, que fundamentam o experimento deste TCC. Esta síntese incorpora os resultados de uma busca complementar dedicada, realizada em julho de 2026 com mais de dez consultas distintas em bases revisadas por pares (ACM DL, IEEE Xplore, Springer, USENIX, NDSS) e complementadas por buscas gerais na web, especificamente voltada a verificar se alguma das lacunas abaixo já havia sido preenchida por literatura publicada entre o fechamento da revisão original e a data desta revisão.

\textbf{Ausência de comparação direta entre SAST tradicional, ferramentas-agente empresariais e híbrido em um mesmo protocolo.} Os estudos comparativos opõem SAST a LLMs \cite{Ozturk2023AIvsSAST, Tamberg2025LLMHarness, Gnieciak2025LLMvsSAST, Zhou2024SASTLLMRepo} ou avaliam híbridos contra suas linhas de base \cite{Ciavatta2026LLMSAST, Ameen2026QASecClaw}, mas nenhum trabalho identificado , nem mesmo entre os localizados na busca complementar de 2026 , compara as três abordagens simultaneamente, sob o mesmo conjunto de dados e as mesmas métricas.

\textbf{Sub-representação do ecossistema Python.} Um mapeamento com 180 estudos de predição de vulnerabilidades encontrou apenas 2 dedicados a Python, contra 86 para C/C++ \cite{Kalouptsoglou2023VulnPrediction}; em revisão com 208 estudos sobre LLMs, Python representa 7,7\% \cite{Germano2025LLMSLR}. A busca complementar de 2026 confirma que a lacuna está diminuindo, mas não fechando: os trabalhos localizados com foco em Python continuam sendo, quase sem exceção, preprints \cite{Pellew2026RealVuln, Marquer2026SAGA, Tsigkourakos2026QRS, Zhou2024SASTLLMRepo, Adnan2026ALPHA} ou usam amostras muito reduzidas \cite{Farasat2025PythonSecurity}; a única exceção revisada por pares com dados reais de Python em escala (Fu et al., ACM TOSEM) trata de geração e correção de código, não de detecção de vulnerabilidades pré-existentes \cite{Fu2025CopilotSecurity}.

\textbf{Ferramentas empresariais de desenvolvimento assistido por IA não avaliadas como detectores.} Esta é a lacuna mais crítica identificada neste levantamento. Não foi localizado estudo revisado por pares que avalie Cursor ou OpenAI Codex, configurados como ferramentas-agente de desenvolvimento, para detecção de vulnerabilidades em aplicações Python. As evidências mais próximas são agentes de pesquisa construídos com LLMs via API \cite{Yildiz2025LLMAgents, Widyasari2026VulTrial}, um protótipo de detecção e reparo integrado ao IDE \cite{Steenhoek2025IDEUserStudy}, revisão de código por GitHub Copilot \cite{Amro2025CopilotReview} e avaliação de Cursor para geração segura de código \cite{Nie2025SecCodePLT}. Esses estudos reforçam a relevância da comparação, mas não substituem a avaliação dos dois produtos-alvo como detectores de vulnerabilidades existentes em repositórios Python.

\textbf{Reprodutibilidade, não-determinismo e dependência de prompt.} LLMs variam entre execuções idênticas \cite{Ozturk2023AIvsSAST}, seus resultados são fortemente sensíveis à estratégia de \textit{prompt} com ganhos de até 0,18 em F1 entre estratégias \cite{Khare2025LLMEffectiveness} e diferenças superiores a 0,3 em F1 entre formulações de instrução \cite{Iranmanesh2025ZeroFalse} , e estratégias supostamente superiores, como \textit{few-shot}, podem degradar o desempenho \cite{Lin2025LLMCompare}. Protocolos padronizados de avaliação, como \textit{harnesses} dedicados, são resposta ainda incipiente \cite{Tamberg2025LLMHarness}.

\textbf{Fragilidade dos benchmarks.} Vazamento de dados, rótulos incorretos e distribuições irreais inflam os resultados reportados \cite{Ding2025CodeLMVuln, Chen2023DiverseVul, Chakraborty2022DLLimits}; avaliações em nível de função superestimam a capacidade em relação ao nível de repositório \cite{Wen2026VulEval}. A própria dispersão de resultados observada na busca complementar de 2026 , por exemplo, a faixa de 90–100\% de detecção por LLM reportada em \cite{Zhou2024SASTLLMRepo}, atipicamente alta frente aos 62,8\% médios de \cite{Khare2025LLMEffectiveness} ilustra esse risco e reforça a necessidade, no experimento do TCC, de conjuntos de dados verificados e métricas resistentes (avaliação em código seguro, penalização de falsos positivos).

\textbf{Custo e limites operacionais das ferramentas de IA.} As avaliações que reportam custo indicam despesas e tempos ordens de magnitude superiores aos do SAST \cite{Tamberg2025LLMHarness, Li2026ProjectScaleLLM}, e o VulTrial acrescenta um dado ponto de referência recente (US\$ 7,46 a US\$ 23,90 por execução completa de um conjunto de teste de 435 pares, a depender do modelo e da arquitetura de agentes) \cite{Widyasari2026VulTrial}. Ainda assim, essa dimensão permanece rara nas comparações, e a busca complementar de 2026 não localizou nenhum estudo acadêmico revisado por pares que reporte custo ou limites de uso (\textit{rate limits}, cotas) especificamente para Cursor, OpenAI Codex ou Claude Code aplicados a detecção de vulnerabilidades apenas material de mercado (documentação de precificação, relatos de usuários), inadequado como evidência acadêmica e por isso não incorporado a este capítulo.

\textbf{Pouca avaliação da utilidade prática dos relatórios.} Apenas uma pequena parte dos estudos abordam explicação \cite{Germano2025LLMSLR}; a indústria demanda diagnósticos acionáveis \cite{Li2025IndustrySAST}; e a utilidade percebida em campo é baixa mesmo quando as métricas de benchmark são razoáveis \cite{Steenhoek2025IDEUserStudy}. Rubricas de corretude, clareza e acionabilidade existem \cite{Mao2025ExplainableVD}, mas raramente são aplicadas em comparações entre abordagens.

Em conjunto, essas lacunas justificam o desenho deste trabalho: uma comparação empírica, sob protocolo único, entre SAST tradicional (Bandit e Semgrep Community Edition), Cursor e OpenAI Codex como ferramentas-agente empresariais, e uma abordagem híbrida SAST + IA, aplicada a código Python, combinando métricas quantitativas (precisão, \textit{recall}, F1, taxa de falsos positivos, tempo e custo) e avaliação qualitativa da clareza, da explicação e da priorização dos relatórios produzidos. A busca complementar realizada especificamente para esta revisão reforça a urgência desse desenho: apesar da expansão de agentes de codificação comerciais e de estudos adjacentes de segurança, não foi localizado estudo revisado por pares que avalie Cursor ou OpenAI Codex como detectores de vulnerabilidades em aplicações Python , o espaço específico que este TCC ocupa.
